\chapter{Appendix: Implementation Details}
\label{app:implementation}

% ============================================================================
% Appendix with additional implementation details, code snippets, etc.
% ============================================================================

\section{Configuration File Example}
\label{app:config_example}

Example YAML configuration for experiments:

\begin{lstlisting}[language=yaml]
# Example experiment configuration
model_type: xgboost
random_state: 42

feature_selection:
  correlation_threshold: 0.95
  top_k_features: 50
  selection_scope: global

model_params:
  max_depth: 6
  learning_rate: 0.1
  n_estimators: 100
  subsample: 0.8

data:
  dataset: BOAS
  n_subjects: 128
  channels: [C3-A2, C4-A1, F3-A2, F4-A1, O1-A2, O2-A1]
  epoch_duration: 30  # seconds
\end{lstlisting}

\section{Feature List}
\label{app:feature_list}

Complete list of 149 extracted features:

% TODO: Auto-generate this from feature_extractor.py
% For each channel (6 total), for each feature type:

\subsection{Time-Domain Features (per channel)}
\label{app:time_features}

\begin{itemize}
    \item Mean
    \item Variance
    \item Standard deviation
    \item Skewness
    \item Kurtosis
    \item Zero-crossing rate
    \item Line length
    \item Hjorth activity
    \item Hjorth mobility
    \item Hjorth complexity
    \item Min/Max amplitude
    \item Peak-to-peak amplitude
\end{itemize}

\subsection{Frequency-Domain Features (per channel)}
\label{app:freq_features}

\begin{itemize}
    \item Total power
    \item Delta band power (0.5--4 Hz)
    \item Theta band power (4--8 Hz)
    \item Alpha band power (8--13 Hz)
    \item Beta band power (13--30 Hz)
    \item Relative delta power
    \item Relative theta power
    \item Relative alpha power
    \item Relative beta power
    \item Delta/Alpha ratio
    \item Theta/Beta ratio
    \item Spectral entropy
    \item Spectral edge frequency (50\%, 90\%, 95\%)
\end{itemize}

\section{Hyperparameter Values}
\label{app:hyperparameters}

% TODO: Fill in from config.py or model params

\subsection{XGBoost Hyperparameters}
\label{app:xgb_params}

\begin{table}[htbp]
    \centering
    \caption{XGBoost Hyperparameter Configuration}
    \begin{tabular}{ll}
        \toprule
        \textbf{Parameter} & \textbf{Value} \\
        \midrule
        max\_depth & [VALUE] \\
        learning\_rate & [VALUE] \\
        n\_estimators & [VALUE] \\
        subsample & [VALUE] \\
        colsample\_bytree & [VALUE] \\
        min\_child\_weight & [VALUE] \\
        gamma & [VALUE] \\
        reg\_alpha & [VALUE] \\
        reg\_lambda & [VALUE] \\
        random\_state & 42 \\
        \bottomrule
    \end{tabular}
\end{table}

\subsection{Random Forest Hyperparameters}
\label{app:rf_params}

\begin{table}[htbp]
    \centering
    \caption{Random Forest Hyperparameter Configuration}
    \begin{tabular}{ll}
        \toprule
        \textbf{Parameter} & \textbf{Value} \\
        \midrule
        n\_estimators & [VALUE] \\
        max\_depth & [VALUE] \\
        min\_samples\_split & [VALUE] \\
        min\_samples\_leaf & [VALUE] \\
        max\_features & [VALUE] \\
        bootstrap & [VALUE] \\
        random\_state & 42 \\
        n\_jobs & -1 \\
        \bottomrule
    \end{tabular}
\end{table}

\section{System Requirements}
\label{app:system_requirements}

\subsection{Hardware}
\label{app:hardware}

Minimum hardware requirements for running the pipeline:

\begin{itemize}
    \item \textbf{CPU:} Multi-core processor (8+ cores recommended for parallelization)
    \item \textbf{RAM:} 16 GB minimum, 32 GB recommended
    \item \textbf{Storage:} 50 GB free space for dataset and cache
    \item \textbf{GPU:} Optional (XGBoost GPU support via CUDA)
\end{itemize}

\subsection{Software}
\label{app:software}

Python dependencies (\texttt{requirements.txt}):

\begin{lstlisting}
numpy>=1.21.0
pandas>=1.3.0
scipy>=1.7.0
scikit-learn>=1.0.0
xgboost>=1.5.0
matplotlib>=3.4.0
seaborn>=0.11.0
joblib>=1.0.0
tqdm>=4.62.0
pyyaml>=5.4.0
\end{lstlisting}

% TODO: Add actual versions from requirements.txt

\section{Cache File Structure}
\label{app:cache_structure}

\subsection{Feature Cache Directory Structure}
\label{app:feature_cache_structure}

\begin{lstlisting}
results/
  features_cache_global/
    subject_001/
      features.npy
      labels.npy
      metadata.json
      .fingerprint
    subject_002/
      ...
    ...
    subject_128/
      ...
\end{lstlisting}

\subsection{LOSO Model Cache Directory Structure}
\label{app:loso_cache_structure}

\begin{lstlisting}
results/
  loso_cache/
    <fingerprint_1>/
      subject_001/
        model.pkl
        predictions.npy
        metadata.json
      subject_002/
        ...
      ...
    <fingerprint_2>/
      ...
\end{lstlisting}

\section{Execution Time Breakdown}
\label{app:execution_time}

% TODO: Fill in after running experiments

Approximate execution times per stage (on [HARDWARE SPEC]):

\begin{table}[htbp]
    \centering
    \caption{Pipeline Execution Time Breakdown}
    \begin{tabular}{lrr}
        \toprule
        \textbf{Stage} & \textbf{Time (first run)} & \textbf{Time (cached)} \\
        \midrule
        Data Loading & X min & X min \\
        Feature Extraction (128 subj) & X hours & <1 min \\
        Feature Selection (per fold) & X sec & X sec \\
        Model Training (per fold) & X min & <1 sec \\
        \midrule
        \textbf{Total (full grid)} & X hours & X hours \\
        \bottomrule
    \end{tabular}
\end{table}

\section{Code Repository}
\label{app:code_repository}

The complete implementation is available at:

% TODO: Add GitHub URL or archive link
\begin{itemize}
    \item \textbf{Repository:} [URL]
    \item \textbf{License:} [LICENSE TYPE]
    \item \textbf{Documentation:} See \texttt{README.md} and \texttt{docs/} directory
\end{itemize}

\section{Reproducing Results}
\label{app:reproduction}

To reproduce the experiments:

\begin{enumerate}
    \item Clone the repository
    \item Install dependencies: \texttt{pip install -r requirements.txt}
    \item Download BOAS dataset to \texttt{data/boas/}
    \item Run feature extraction: \texttt{python run\_experiment.py --full}
    \item Run training: \texttt{python run\_training.py --full}
    \item Results will be saved to \texttt{results/}
\end{enumerate}

% ============================================================================
% Add more appendices as needed (B, C, etc.)
% ============================================================================
